\clearpage
\section{Holonomy-twisted lane transport and orbit descent}
\label{sec:block16-holonomy-twisted-lane-transport}

The finite atlas theorem in \cref{thm:block15-lane-descent} imposed trivial
address holonomy in order to obtain one globally fixed set of lane labels.  That
hypothesis is sufficient for ordinary globally addressed lanes, but it is
stronger than the hypothesis needed for transport itself.  When addresses are
permuted around atlas loops, the local lane matrices still glue.  They glue as
sections of a finite matrix local system rather than as one matrix with a fixed
global row and column labeling.

This section isolates that distinction.  The local analysis, synthesis,
composition, refinement, selector, and residual laws survive arbitrary finite
address holonomy.  The holonomy orbits then give the finest ordinary global
coarsening obtainable by summing the original lane projectors.  Thus nontrivial
holonomy is an obstruction to untwisted labels, not an obstruction to the
finite transport calculus.

\subsection{Holonomy-profiled lane data}
\label{subsec:block16-holonomy-profiled-data}

Let \(\mathscr G\) be a finite connected groupoid of charts.  For an arrow
\(\gamma:U\to V\), let
\[
 F_\gamma:\mathcal C_U\longrightarrow\mathcal C_V
\]
be an additive equivalence between additive idempotent-complete categories.
The transition functors carry specified identity and composition comparisons.
To keep the formulas readable, those comparison isomorphisms are suppressed;
every displayed cocycle equation is understood after inserting the declared
comparisons.

A descent object
\[
 A=(A_U,\theta^A_\gamma)
\]
consists of objects \(A_U\in\mathcal C_U\) and isomorphisms
\[
 \theta^A_\gamma:F_\gamma(A_U)\xrightarrow{\sim}A_V
\]
satisfying the identity and cocycle equations.  A descended morphism
\(f:X\to Y\) is a compatible family \((f_U)_U\) satisfying
\[
 \theta^Y_\gamma F_\gamma(f_U)=f_V\theta^X_\gamma.
\]

\begin{definition}[Holonomy-profiled finite lane datum]
\label{def:block16-holonomy-profiled-lane-datum}
Fix a finite resolution shape \(\mathsf R\).  A holonomy-profiled finite lane
datum on a descent object \(A\) consists of the following information.

For every chart \(U\) and resolution \(r\in\mathsf R\), the object \(A_U\)
has a finite profiled lane system
\[
 \left(I_{A,U,r},p^{r}_{A,U,i},A^r_{U,i},
 \iota^r_{A,U,i},\pi^r_{A,U,i}\right)_{i\in I_{A,U,r}}
\]
in the sense of \cref{def:block15-finite-profiled-lane-system}.  For every
arrow \(\gamma:U\to V\), there is an address bijection
\[
 \lambda^A_{\gamma,r}:I_{A,U,r}\xrightarrow{\sim}I_{A,V,r}
\]
and, for every address \(i\), a splitting comparison
\[
 \beta^A_{\gamma,r,i}:
 F_\gamma(A^r_{U,i})\xrightarrow{\sim}
 A^r_{V,\lambda^A_{\gamma,r}(i)}.
\]
They satisfy the address and splitting cocycle laws and the two compatibility
identities
\[
 \theta^A_\gamma F_\gamma(\iota^r_{A,U,i})
 =
 \iota^r_{A,V,\lambda^A_{\gamma,r}(i)}
 \beta^A_{\gamma,r,i},
\]
\[
 \beta^A_{\gamma,r,i}F_\gamma(\pi^r_{A,U,i})
 =
 \pi^r_{A,V,\lambda^A_{\gamma,r}(i)}\theta^A_\gamma.
\]
For \(s\succeq r\), the refinement maps are equivariant:
\[
 \rho^{A,V}_{sr}\lambda^A_{\gamma,s}
 =
 \lambda^A_{\gamma,r}\rho^{A,U}_{sr}.
\]
No global address set and no global trivialization of the address bijections is
assumed.
\end{definition}

\begin{definition}[Matrix local system]
\label{def:block16-matrix-local-system}
Let \(X,Y\) carry holonomy-profiled lane data.  At source resolution \(r\)
and target resolution \(u\), define the matrix group over a chart \(U\) by
\[
 \mathscr M^{u,r}_{Y,X}(U)
 :=
 \bigoplus_{a\in I_{Y,U,u}}
 \bigoplus_{i\in I_{X,U,r}}
 \Hom_{\mathcal C_U}(X^r_{U,i},Y^u_{U,a}).
\]
For \(h=(h_{ai})\in\mathscr M^{u,r}_{Y,X}(U)\), define transport along
\(\gamma:U\to V\) by
\[
 \left(\gamma_*h\right)_{
 \lambda^Y_{\gamma,u}(a),\lambda^X_{\gamma,r}(i)}
 :=
 \beta^Y_{\gamma,u,a}
 F_\gamma(h_{ai})
 \left(\beta^X_{\gamma,r,i}\right)^{-1}.
\]
The group of compatible sections is
\[
 \Gamma\!\left(\mathscr G,\mathscr M^{u,r}_{Y,X}\right)
 :=
 \left\{(h_U)_U:\gamma_*h_U=h_V
 \text{ for every }\gamma:U\to V\right\}.
\]
\end{definition}

\subsection{The holonomy-twisted transport theorem}
\label{subsec:block16-main-theorem}

\begin{theorem}[Holonomy-Twisted Universal Lane Transport Theorem]
\label{thm:block16-holonomy-twisted-lane-transport}
Let \(X,Y,Z\) be descent objects over \(\mathscr G\), each equipped with a
holonomy-profiled finite lane datum over the same finite resolution shape.
Then the following statements hold.

\begin{enumerate}[label=\textup{(H\arabic*)}]
\item \textbf{Matrix local system.}
The maps \(\gamma_*\) of
\cref{def:block16-matrix-local-system} are additive isomorphisms and satisfy
the identity and cocycle laws.  Hence
\(\mathscr M^{u,r}_{Y,X}\) is a local system of abelian groups with finitely many matrix positions in every chart.

\item \textbf{Exact twisted analysis and synthesis.}
Lane analysis induces an isomorphism
\[
 \mathfrak A^{u,r}_{Y,X}:
 \Hom_{\Desc(\mathscr G)}(X,Y)
 \xrightarrow{\sim}
 \Gamma\!\left(\mathscr G,\mathscr M^{u,r}_{Y,X}\right),
 \qquad
 f\longmapsto
 \left(\bigl([f_U]^{u,r}_{ai}\bigr)_{a,i}\right)_U.
\]
Its inverse is chartwise lane synthesis:
\[
 (h_U)_U\longmapsto
 \left(
 \sum_{a,i}
 \iota^u_{Y,U,a}h_{U,ai}\pi^r_{X,U,i}
 \right)_U.
\]
Thus nontrivial address holonomy does not destroy the complete matrix record.
It changes the global type of that record from a constantly labeled matrix to
a compatible section of a matrix local system.

\item \textbf{Twisted matrix composition.}
For
\[
 h\in\Gamma(\mathscr G,\mathscr M^{u,r}_{Y,X}),
 \qquad
 k\in\Gamma(\mathscr G,\mathscr M^{v,u}_{Z,Y}),
\]
define chartwise contraction by
\[
 (k\star h)_{U,ci}
 :=
 \sum_{a\in I_{Y,U,u}}k_{U,ca}h_{U,ai}.
\]
Then \(k\star h\) is a compatible section of
\(\mathscr M^{v,r}_{Z,X}\), and
\[
 \mathfrak A^{v,r}_{Z,X}(gf)
 =
 \mathfrak A^{v,u}_{Z,Y}(g)
 \star
 \mathfrak A^{u,r}_{Y,X}(f).
\]
The diagonal identity matrices are compatible sections.  More precisely, fix
one declared resolution \(r_A\) for each lane-decorated descent object \(A\).
The category with the same objects, morphism groups
\[
 \Gamma\!\left(\mathscr G,
 \mathscr M^{r_B,r_A}_{B,A}\right),
\]
and composition \(\star\) is isomorphic, by analysis and synthesis, to the
corresponding category of lane-decorated descent objects and their underlying
descended morphisms.

\item \textbf{Equivariant refinement.}
Every certified fine-to-coarse coordinate change from
\cref{prop:block15-certified-coarsening} commutes with the transition maps
\(\gamma_*\).  It therefore defines an isomorphism of matrix local systems.
These isomorphisms compose strictly over chains of declared resolutions and
intertwine twisted matrix contraction.

\item \textbf{Holonomy-compatible selectors and residuals.}
Let
\[
 E_U\subseteq I_{Y,U,u}\times I_{X,U,r}
\]
be a family satisfying
\[
 E_V=
 (\lambda^Y_{\gamma,u}\times\lambda^X_{\gamma,r})(E_U)
\]
for every arrow \(\gamma:U\to V\).  Then the local selectors
\[
 C_{E_U}(h_U)_{ai}
 =
 \begin{cases}
 h_{U,ai},&(a,i)\in E_U,\\
 0,&(a,i)\notin E_U
 \end{cases}
\]
assemble to an idempotent endomorphism \(C_E\) of the section group.  For two
compatible support families \(E,F\),
\[
 C_EC_F=C_{E\cap F},
 \qquad
 R_E:=1-C_E,
 \qquad
 h=C_Eh+R_Eh.
\]
Hence named residuals remain exact under arbitrary finite address holonomy,
provided the retained support itself is holonomy compatible.
\end{enumerate}
\end{theorem}

\begin{proof}
For (H1), additivity and invertibility follow from the corresponding
properties of \(F_\gamma\) and the splitting comparisons \(\beta\).  The
address and splitting cocycle laws give
\((\delta\gamma)_*=\delta_*\gamma_*\) after the declared transition
comparisons.

For (H2), let \(f=(f_U)_U\) be descended.  For a transition
\(\gamma:U\to V\), the transported \((a,i)\)-block is
\[
 \begin{aligned}
 &\beta^Y_{\gamma,u,a}
 F_\gamma\!\left(
 \pi^u_{Y,U,a}f_U\iota^r_{X,U,i}
 \right)
 \left(\beta^X_{\gamma,r,i}\right)^{-1}\\
 &\quad=
 \pi^u_{Y,V,\lambda^Y(a)}
 \theta^Y_\gamma F_\gamma(f_U)(\theta^X_\gamma)^{-1}
 \iota^r_{X,V,\lambda^X(i)}\\
 &\quad=
 \pi^u_{Y,V,\lambda^Y(a)}
 f_V
 \iota^r_{X,V,\lambda^X(i)},
 \end{aligned}
\]
which is the corresponding block of \(f_V\).  Thus analysis produces a
compatible section.

Conversely, let \(h=(h_U)_U\) be a compatible section and synthesize
\[
 f_U:=\sum_{a,i}\iota^u_{Y,U,a}h_{U,ai}\pi^r_{X,U,i}.
\]
Transporting this sum, applying the splitting comparison identities, and
reindexing through the address bijections gives \(f_V\).  Thus \((f_U)_U\) is
a descended morphism.  Chartwise analysis and synthesis are mutually inverse
by \cref{thm:block15-universal-lane-transport}; therefore their restrictions
to compatible families are mutually inverse.

For (H3), the local path law gives the displayed contraction formula in every
chart.  Applying \(\gamma_*\) to a contracted block distributes over the finite
sum and carries each product to the product of the transported blocks.  Since
\(\lambda^Y_{\gamma,u}\) is a bijection, reindexing the intermediate address
shows
\[
 \gamma_*(k_U\star h_U)
 =
 (\gamma_*k_U)\star(\gamma_*h_U)
 =
 k_V\star h_V.
\]
The identity section is compatible by the splitting cocycle.  The equivalence
then follows from (H2).

For (H4), the local fine-to-coarse coordinate map is analysis after synthesis.
Both operations commute with transition by (H2), so their composite does as
well.  The strict composition and contraction identities are the corresponding
chartwise identities from
\cref{cor:block15-functorial-coarsening,thm:block15-matrix-category-equivalence}.

For (H5), support compatibility implies
\(\gamma_*C_{E_U}=C_{E_V}\gamma_*\), so the selectors preserve compatible
sections.  Every Boolean and residual identity is then checked chartwise.
\end{proof}


\begin{corollary}[Base-fiber invariant normal form]
\label{cor:block16-base-fiber-invariant-normal-form}
Fix a base chart \(U_0\).  Evaluation at \(U_0\) induces an isomorphism
\[
 \Gamma\!\left(\mathscr G,\mathscr M^{u,r}_{Y,X}\right)
 \xrightarrow{\sim}
 \mathscr M^{u,r}_{Y,X}(U_0)^{\Aut_{\mathscr G}(U_0)},
\]
where the right-hand side is the subgroup fixed by every loop-transport map.
Consequently,
\[
 \Hom_{\Desc(\mathscr G)}(X,Y)
 \cong
 \mathscr M^{u,r}_{Y,X}(U_0)^{\Aut_{\mathscr G}(U_0)}.
\]
After choosing one arrow \(\tau_U:U_0\to U\) for each chart, the inverse sends
a fixed base matrix \(h_0\) to the section
\(h_U=(\tau_U)_*h_0\).  The result is independent of the chosen arrows because
two choices differ by loop transport at \(U_0\).
\end{corollary}

\begin{proof}
A compatible section evaluates to a loop-fixed element at \(U_0\).  Conversely,
for an arrow \(\gamma:U\to V\), the two arrows
\(\gamma\tau_U\) and \(\tau_V\) differ by a loop at \(U_0\); loop invariance
therefore gives \(\gamma_*h_U=h_V\).  Evaluation and this reconstruction are
mutually inverse.  Compose with the analysis isomorphism in (H2).
\end{proof}

\begin{remark}[What has and has not been removed]
\label{rem:block16-removed-triviality-only}
The theorem removes only the global address trivializations
\(\mu_U^{A,r}\) from \cref{thm:block15-lane-descent}.  It does not remove the
address cocycle, refinement naturality, projector compatibility, splitting
comparisons, or additive preservation hypotheses.  Without those data, even
the matrix transition maps are undefined or fail their cocycle law.
\end{remark}

\subsection{Monodromy orbits and the finest untwisted coarsening}
\label{subsec:block16-orbit-descent}

Fix a base chart \(U_0\).  Every loop \(\gamma\in\Aut_{\mathscr G}(U_0)\)
acts on an address set by \(i\mapsto\lambda^A_{\gamma,r}(i)\).  Write
\[
 H_A^r
 :=
 \operatorname{im}\!\left(
 \Aut_{\mathscr G}(U_0)
 \longrightarrow
 \operatorname{Sym}(I_{A,U_0,r})
 \right)
\]
for the resulting finite monodromy group.  For a source-target pair, the same
loop acts diagonally by
\[
 (a,i)\longmapsto
 \left(\lambda^Y_{\gamma,u}(a),
 \lambda^X_{\gamma,r}(i)\right).
\]
Its image in the symmetric group of
\(I_{Y,U_0,u}\times I_{X,U_0,r}\) is denoted
\(H_{Y,X}^{u,r}\).  This is the simultaneous monodromy image; it need not be
the direct product of the two separate monodromy groups.

\begin{theorem}[Orbit descent and the finest globally labeled coarsening]
\label{thm:block16-orbit-descent}
Under the hypotheses of
\cref{thm:block16-holonomy-twisted-lane-transport}, the following statements
hold.

\begin{enumerate}[label=\textup{(O\arabic*)}]
\item For every orbit \(\mathcal O\subseteq I_{A,U_0,r}\) of \(H_A^r\), path
transport defines a path-independent orbit subset
\(\mathcal O_U\subseteq I_{A,U,r}\) in every chart.  The local projector sums
\[
 P^r_{A,U,\mathcal O}
 :=
 \sum_{i\in\mathcal O_U}p^r_{A,U,i}
\]
descend to a split idempotent on \(A\).  As \(\mathcal O\) ranges over the
monodromy orbits, these projectors are pairwise annihilating and sum to
\(1_A\).

\item The orbit projectors form the finest ordinary globally labeled
coarsening whose local projectors are sums of the original lane projectors.
More precisely, any globally labeled projector of the form
\[
 Q_U=\sum_{i\in S_U}p^r_{A,U,i}
\]
with \(S_V=\lambda^A_{\gamma,r}(S_U)\) and with the same global label preserved
around every loop has \(S_{U_0}\) equal to a union of \(H_A^r\)-orbits.

\item A particular local lane descends as an individually labeled global lane
exactly when its address is fixed by \(H_A^r\).  All original lane labels
descend simultaneously exactly when the monodromy action on the address set is
trivial.

\item Holonomy-compatible support families in (H5) are in bijection with
unions of \(H_{Y,X}^{u,r}\)-orbits in the base pair-address set.  The minimal
nonzero compatible Boolean selectors are therefore the selectors associated
with single pair orbits.

\item If \(s\succeq r\), equivariance of the refinement map sends every fine
monodromy orbit onto a coarse monodromy orbit.  For a coarse orbit
\(\mathcal O\),
\[
 P^r_{A,\mathcal O}
 =
 \sum_{\substack{\mathcal Q\text{ a fine orbit}\\
 \mathcal Q\subseteq(\rho^A_{sr})^{-1}(\mathcal O)}}
 P^s_{A,\mathcal Q}.
\]
Thus orbit descent and projector-sum refinement commute exactly.
\end{enumerate}
\end{theorem}

\begin{proof}
Choose a path from \(U_0\) to \(U\) and transport \(\mathcal O\) along its
address bijection.  Two path choices differ by a loop at \(U_0\), which
preserves \(\mathcal O\) setwise.  Hence \(\mathcal O_U\) is path independent.
Projector compatibility gives
\[
 \theta^A_\gamma F_\gamma(P^r_{A,U,\mathcal O})
 =
 P^r_{A,V,\mathcal O}\theta^A_\gamma,
\]
so the orbit sums descend.  Their local splittings are the finite direct sums
\(\bigoplus_{i\in\mathcal O_U}A^r_{U,i}\); the direct sums of the \(\beta\)-maps,
with the induced permutation of summands, provide their descent data.  This
proves (O1).

For (O2), a globally fixed label must return to itself after every loop, so
\(S_{U_0}\) is invariant under \(H_A^r\).  Every invariant subset of a finite
group action is a union of orbits.  Therefore every such globally labeled
coarsening is coarser than the orbit partition.  Statement (O3) is the
singleton case.

The same path-independence argument on the pair-address set shows that a
compatible support family is determined by an invariant base subset, and every
invariant base subset is a union of pair orbits.  This proves (O4).

Finally, equivariance of \(\rho^A_{sr}\) implies that the image of a fine orbit
is the orbit of the image of any one of its points.  Hence the preimage of a
coarse orbit is a union of fine orbits.  Summing the local projector-refinement
identity over all addresses in the coarse orbit gives the formula in (O5).
\end{proof}

\begin{remark}[Object orbits and matrix orbits differ]
\label{rem:block16-object-vs-pair-orbits}
The orbit lanes on an object are controlled by \(H_A^r\backslash I_{A,r}\).
Transport selectors are controlled by the generally finer orbit set
\[
 H_{Y,X}^{u,r}\backslash
 \left(I_{Y,u}\times I_{X,r}\right).
\]
Even when an object has only one transitive address orbit, its endomorphism
matrices can have several pair orbits.  Holonomy can therefore erase individual
object labels while retaining a nontrivial global transport algebra.
\end{remark}

\subsection{Permutation modules and the orbital path algebra}
\label{subsec:block16-permutation-orbital-calculus}

The one-chart groupoid with automorphism group \(H\) gives a concrete exact
model.  It also identifies the algebra carried by the pair-orbit selectors.

\begin{theorem}[Permutation-module orbital lane calculus]
\label{thm:block16-permutation-orbital-calculus}
Let \(\kk\) be a commutative ring, let a finite group \(H\) act on finite sets
\(I,J,K\), and let \(\kk[I],\kk[J],\kk[K]\) be the corresponding permutation
modules.  For an orbit
\(\mathcal O\subseteq J\times I\) under the diagonal action, let
\(B_{\mathcal O}:\kk[I]\to\kk[J]\) be the matrix whose \((j,i)\)-entry is one
when \((j,i)\in\mathcal O\) and zero otherwise.  Then
\[
 \Hom_{\kk H}(\kk[I],\kk[J])
 =
 \bigoplus_{\mathcal O\in H\backslash(J\times I)}
 \kk B_{\mathcal O}.
\]
In particular, an equivariant matrix is exactly a matrix constant on every
pair orbit.

Let
\[
 \mathcal P\subseteq J\times I,
 \qquad
 \mathcal Q\subseteq K\times J,
 \qquad
 \mathcal R\subseteq K\times I
\]
be diagonal \(H\)-orbits.  For any \((k,i)\in\mathcal R\), define
\[
 c^{\mathcal R}_{\mathcal Q,\mathcal P}
 :=
 \#\left\{
 j\in J:(k,j)\in\mathcal Q,
 (j,i)\in\mathcal P
 \right\}.
\]
This integer is independent of the chosen \((k,i)\in\mathcal R\), and
\[
 B_{\mathcal Q}B_{\mathcal P}
 =
 \sum_{\mathcal R\in H\backslash(K\times I)}
 c^{\mathcal R}_{\mathcal Q,\mathcal P}B_{\mathcal R}.
\]
Thus the lane path law descends to an exact orbit basis with nonnegative
integer structure constants.
\end{theorem}

\begin{proof}
Write an arbitrary \(\kk\)-linear map as a matrix \(A=(A_{ji})\).  The
intertwining equations for the permutation actions are equivalent to
\[
 A_{h j,h i}=A_{j i}
 \qquad(h\in H).
\]
Hence the entries are constant on diagonal orbits.  The orbit-indicator
matrices have disjoint supports, so they are linearly independent over every
commutative ring and span all invariant matrices.

The \((k,i)\)-entry of
\(B_{\mathcal Q}B_{\mathcal P}\) is the displayed cardinality.  If
\((k',i')=(hk,hi)\), then \(j\mapsto hj\) is a bijection between the two sets
being counted.  The entry is therefore constant on \(\mathcal R\), which gives
the orbit-basis expansion.
\end{proof}

\begin{corollary}[Orbital endomorphism algebra]
\label{cor:block16-orbital-endomorphism-algebra}
For one finite \(H\)-set \(I\), the pair-orbit matrices form a basis of
\(\End_{\kk H}(\kk[I])\).  Matrix multiplication in this basis is exactly the
finite compatible-path contraction of
\cref{thm:block16-holonomy-twisted-lane-transport}, with the preceding
intersection numbers as structure constants.
\end{corollary}

\begin{example}[The smallest nontrivial holonomy]
\label{ex:block16-c2-swap}
Let \(H=C_2=\{1,\sigma\}\) act on \(I=\{1,2\}\) by swapping the two addresses.
On \(X=\kk^2\), let
\[
 S=
 \begin{pmatrix}0&1\\1&0\end{pmatrix},
 \qquad
 p_1=
 \begin{pmatrix}1&0\\0&0\end{pmatrix},
 \qquad
 p_2=
 \begin{pmatrix}0&0\\0&1\end{pmatrix}.
\]
The individual projectors do not descend because
\(Sp_1S^{-1}=p_2\), while their orbit sum \(p_1+p_2=1\) does descend.  The
pair-address set has two orbits,
\[
 \Delta=\{(1,1),(2,2)\},
 \qquad
 \Omega=\{(1,2),(2,1)\}.
\]
Accordingly,
\[
 \End_{\kk C_2}(\kk^2)
 =
 \left\{
 \begin{pmatrix}a&b\\b&a\end{pmatrix}:a,b\in\kk
 \right\}
 =
 \kk B_\Delta\oplus\kk B_\Omega.
\]
The singleton endpoint selector \(\{(1,1)\}\) is not holonomy compatible and
sends the equivariant identity matrix to \(p_1\), which is not equivariant.
The orbit selectors \(C_\Delta\) and \(C_\Omega\) are compatible and split every
equivariant endomorphism into its diagonal-orbit and off-diagonal-orbit parts.
\end{example}

\subsection{Cyclotomic orbit lanes over the rational field}
\label{subsec:block16-cyclotomic-orbit-lanes}

The split cyclic group algebra supplies an arithmetic realization in which the
address monodromy and its orbit coarsening are explicit.

\begin{proposition}[Rational descent of split cyclic Fourier lanes]
\label{prop:block16-rational-cyclotomic-orbit-lanes}
Let \(N\ge1\), let \(C_N=\langle g\rangle\), and let
\(K=\mathbb Q(\zeta_N)\).  In \(K[C_N]\), define
\[
 e_j
 :=
 \frac1N\sum_{\ell=0}^{N-1}
 \zeta_N^{-j\ell}g^\ell,
 \qquad j\in\mathbb Z/N\mathbb Z.
\]
These are the primitive split Fourier idempotents.  For
\(\sigma_a\in\Gal(K/\mathbb Q)\), where
\(\sigma_a(\zeta_N)=\zeta_N^a\),
\[
 \sigma_a(e_j)=e_{aj}.
\]
The Galois orbits are indexed by divisors \(q\mid N\), with
\[
 \mathcal O_q
 =
 \left\{j\in\mathbb Z/N\mathbb Z:
 \frac{N}{\gcd(N,j)}=q\right\}.
\]
Their orbit sums
\[
 E_q:=\sum_{j\in\mathcal O_q}e_j
\]
belong to \(\mathbb Q[C_N]\) and satisfy
\[
 E_qE_{q'}=\delta_{q,q'}E_q,
 \qquad
 \sum_{q\mid N}E_q=1.
\]
More explicitly,
\[
 \boxed{
 E_q
 =
 \frac1N\sum_{\ell=0}^{N-1}c_q(\ell)g^\ell,}
\]
where
\[
 c_q(\ell)
 :=
 \sum_{\substack{1\le a\le q\\\gcd(a,q)=1}}
 e^{2\pi i a\ell/q}
 =
 \sum_{d\mid\gcd(q,\ell)}d\,\mu(q/d)
 \in\mathbb Z
\]
is the Ramanujan sum.  The projectors \((E_q)_{q\mid N}\) form the finest
\(\mathbb Q\)-rational coarsening obtainable by summing the split Fourier
idempotents, and
\[
 \mathbb Q[C_N]E_q\cong\mathbb Q(\zeta_q).
\]
\end{proposition}

\begin{proof}
Character orthogonality gives the primitive idempotent relations for the
\(e_j\).  Applying \(\sigma_a\) to their coefficients gives
\(\sigma_a(e_j)=e_{aj}\).  The unit group
\((\mathbb Z/N\mathbb Z)^\times\) acts transitively on elements of
\(\mathbb Z/N\mathbb Z\) having the same additive order, so its orbits are the
sets \(\mathcal O_q\).

Write every \(j\in\mathcal O_q\) uniquely as
\(j=(N/q)a\) with \(a\in(\mathbb Z/q\mathbb Z)^\times\).  The coefficient of
\(g^\ell\) in \(E_q\) is then
\[
 \frac1N
 \sum_{a\in(\mathbb Z/q\mathbb Z)^\times}
 \zeta_q^{-a\ell}
 =
 \frac{c_q(\ell)}{N},
\]
where the last equality uses the integrality and evenness of the Ramanujan
sum.  Hence \(E_q\in\mathbb Q[C_N]\).  Orthogonality and completeness follow
by grouping the primitive \(e_j\) by their disjoint Galois orbits.  Any rational
sum of split primitive idempotents must be Galois invariant, so its support is a
union of these orbits; this proves fineness.  Finally,
\[
 \mathbb Q[C_N]
 \cong
 \mathbb Q[x]/(x^N-1)
 \cong
 \prod_{q\mid N}\mathbb Q[x]/(\Phi_q(x))
 \cong
 \prod_{q\mid N}\mathbb Q(\zeta_q),
\]
and \(E_q\) is the identity projector of the \(q\)-th factor.
\end{proof}

\begin{boundary}[Cyclic character lanes are not coefficient-residue lanes]
\label{bnd:block16-cyclic-character-residue-distinction}
The preceding proposition concerns primitive character idempotents in the
split group algebra \(K[C_N]\) and their rational Galois descent.  The
coefficient-residue projectors \(\mathcal P_{N,j}\) on formal series from
\cref{thm:block15-factorial-residue-realization} are defined coefficientwise
over every commutative ring and already have fixed residue labels.  A chosen
rotation representation relates the two constructions after scalar extension,
but their descent data are not identical and should not be conflated.  Nothing
in orbit descent turns the factorial step-product into an additive idempotent
splitting.
\end{boundary}

\subsection{Exact finite validation and scope}
\label{subsec:block16-validation-scope}

The accompanying verifier tests finite permutation actions, pair-orbit bases,
equivariant matrix composition, orbit-compatible and deliberately incompatible
selectors, equivariant refinement, and the exact Ramanujan-sum projectors in
\(\mathbb Q[C_N]\).  It uses exact integer arithmetic after clearing the common cyclotomic
denominators.  These tests are
an implementation regression for the displayed formulas, not a substitute for
the proofs and not an external validation record.

\IfFileExists{generated/block16_validation_table.tex}{%
% Generated by verify_block16_holonomy_lane_transport.py
\begin{table}[htbp]
\centering
\caption{Exact finite regression for selected holonomy-twisted lane clauses. The table is an implementation record, not a proof or external validation.}
\label{tab:block16-validation}
\begin{tabular}{@{}lr@{}}
\toprule
Validation family & Exact assertions\\
\midrule
Group actions and object-orbit descent & 7,780\\
Pair-orbit partitions and intertwiner bases & 3,537\\
Orbital path composition and structure constants & 1,212,963\\
Compatible selectors and negative controls & 994\\
Equivariant refinement and orbit-projector sums & 10,162\\
Rational cyclotomic orbit projectors & 109,374\\
Theorem-source contract & 13\\
\midrule
Total selected implementation assertions & 1,344,823\\
Recorded failures & 0\\
\bottomrule
\end{tabular}
\end{table}
%
}{%
\begin{table}[htbp]
\centering
\caption{Exact finite regression for selected holonomy-twisted lane clauses.
Run the Block 16 verifier to generate the complete table.}
\label{tab:block16-validation-placeholder}
\begin{tabular}{@{}ll@{}}
\toprule
Record & Status\\
\midrule
Machine-readable validation & Run \texttt{verify\_block16\_holonomy\_lane\_transport.py}\\
\bottomrule
\end{tabular}
\end{table}%
}

\begin{resultbox}
\textbf{Holonomy-twisted lane transport result.}
Finite nontrivial address holonomy does not invalidate exact lane transport.
The complete global coordinate object is a matrix local system, and descended
morphisms are exactly its compatible sections.  Composition remains finite
path contraction, complete resolution changes remain isomorphisms, and every
holonomy-compatible selector retains an exact named residual.  Monodromy
orbits give the finest ordinary globally labeled coarsening of object lanes,
while simultaneous pair orbits give the atomic global transport selectors.  In
permutation modules these pair orbits form an exact basis of the intertwiner
space with integer path-count structure constants.
\end{resultbox}

\begin{boundary}[Finite twisted scope]
\label{bnd:block16-finite-twisted-scope}
The theorem remains finite and profiled.  It assumes invertible chart
transitions, finite address sets, coherent address and splitting cocycles, and
profile preservation of the operations used.  It does not prove convergence
for countably infinite lane sums, descent across noninvertible bridges, higher
stack effectivity, or a canonical orbit system independent of the initially
declared lane presentation.  Orbit count is not identified with dimension,
rank, entropy, information, or complexity.  The permutation-module and
cyclotomic formulas have classical algebraic antecedents; their role here is to
provide exact native realizations of the twisted lane doctrine.
\end{boundary}
